% ============================================================
%  CHAPTER 23: CURRENT ELECTRICITY
%  The Physics Codex: A Complete University Foundation
%  Korede Samuel Omotosho (Falcon)
%  Aurikrex Learning Systems, 2026
%
%  File: chapters/part4/ch23_current_electricity.tex
% ============================================================

\chapter{Current Electricity}
\label{ch:current_electricity}
\minitoc
\newpage

\chapterquote{The science of today is the technology
of tomorrow.}{Edward Teller}

\begin{objectivebox}
By the end of this chapter, you should be able to:
\begin{enumerate}[noitemsep]
  \item Define electric current, potential difference
  and resistance
  \item State and apply Ohm's Law
  \item Define resistivity and calculate resistance
  from geometry
  \item Describe how resistance varies with temperature
  for metals, thermistors and semiconductors
  \item Combine resistors in series and parallel
  \item Apply Kirchhoff's laws to complex circuits
  \item Define EMF and internal resistance and apply
  them to circuit problems
  \item Define and apply the potential divider
  \item Describe the Wheatstone bridge and its use
  \item Apply the potentiometer principle
\end{enumerate}
\end{objectivebox}

% ============================================================
\section{Intuition: What is Electric Current?}
% ============================================================

Imagine water flowing through a pipe. The rate of
flow --- litres per second --- tells you how much
water is moving. Electric current is exactly analogous:
it is the rate of flow of electric charge through a
conductor.

In a metal wire, the ``water'' is the sea of free
electrons. The electrons drift slowly in one direction
(from negative to positive terminal, though by
convention current flows the other way), colliding
repeatedly with the metal ions as they move. The
\textit{drift velocity} of electrons in a wire is
surprisingly slow --- typically a fraction of a
millimetre per second. But there are so many electrons
that even this slow drift constitutes a very significant
rate of charge flow.

The reason a light switch works instantly --- far faster
than the electrons could physically travel from the
switch to the bulb --- is that the electric field
propagates almost at the speed of light through the
wire. Every electron in the circuit starts moving
almost simultaneously.

% ============================================================
\section{Electric Current}
% ============================================================

\begin{definitionbox}[Electric Current]
\textbf{Electric current} $I$ is the rate of flow
of electric charge:
\[
  I = \frac{\Delta Q}{\Delta t}
  \implies Q = It
\]
SI unit: ampere (A) where
$1\ \text{A} = 1\ \text{C\,s}^{-1}$.

\textbf{Conventional current} flows from positive
to negative terminal (outside the source). This is
opposite to the direction of electron flow.

\textbf{Drift velocity} $v_d$ of charge carriers:
\[
  I = nqAv_d
\]
where $n$ = number density of charge carriers
(m$^{-3}$), $q$ = charge per carrier (C), $A$ =
cross-sectional area of conductor (m$^2$),
$v_d$ = drift velocity (m\,s$^{-1}$).
\end{definitionbox}

\begin{examplebox}[23.1]
\stepcounter{workedexample}
A copper wire carries a current of $3.0\ \text{A}$.
The wire has a cross-sectional area of
$1.5\ \text{mm}^2$ and copper has a free electron
density of $8.5\times10^{28}\ \text{m}^{-3}$.\\
(a) Calculate the drift velocity of electrons.\\
(b) How long does it take for one electron to travel
$1.0\ \text{m}$ along the wire?
\end{examplebox}

\begin{solutionbox}
\textbf{(a)} $A = 1.5\times10^{-6}\ \text{m}^2$,
$q = 1.6\times10^{-19}\ \text{C}$:
\[
  v_d = \frac{I}{nqA}
  = \frac{3.0}{8.5\times10^{28}
    \times 1.6\times10^{-19}
    \times 1.5\times10^{-6}}
  = \frac{3.0}{2.04\times10^4}
  = 1.47\times10^{-4}\ \text{m\,s}^{-1}
\]
About $0.15\ \text{mm\,s}^{-1}$ --- slower than a
snail!

\textbf{(b)} Time to travel $1.0\ \text{m}$:
\[
  t = \frac{d}{v_d}
  = \frac{1.0}{1.47\times10^{-4}}
  = 6800\ \text{s} \approx 113\ \text{min}
\]
Almost two hours for an electron to drift $1\ \text{m}$.
This is why understanding drift velocity versus signal
propagation speed is crucial in electronics.
\end{solutionbox}

% ============================================================
\section{Ohm's Law and Resistance}
% ============================================================

\begin{lawbox}[Ohm's Law]
For a metallic conductor at \textbf{constant
temperature}, the current through the conductor is
directly proportional to the potential difference
across it:
\[
  V \propto I \implies V = IR
\]
where $R$ is the \textbf{resistance} of the
conductor. SI unit: ohm ($\Omega$) where
$1\ \Omega = 1\ \text{V\,A}^{-1}$.

This defines resistance for any component:
\[
  R = \frac{V}{I}
\]
A component that obeys Ohm's Law is called
\textbf{ohmic}. One that does not is \textbf{non-ohmic}.
\end{lawbox}

\begin{diagbox}[$I$-$V$ Curves: Ohmic and Non-ohmic (schematic; ranges illustrative)]
\begin{center}
\begin{tikzpicture}
  % Ohmic resistor
  \begin{axis}[
    name=p1,
    width=4.0cm, height=4.2cm,
    xlabel={$V$ (V)}, ylabel={$I$ (A)},
    xmin=0, xmax=5, ymin=0, ymax=5,
    grid=major, grid style={gray!15},
    xtick={0,1,2,3,4,5},
    ytick={0,1,2,3,4,5},
    title={Ohmic Resistor}
  ]
    \addplot[codexblue, very thick, domain=0:5]{x};
  \end{axis}

  % Filament bulb
  \begin{axis}[
    name=p2,
    at={(p1.right of south east)},
    anchor=left of south west,
    xshift=0.25cm,
    width=4.0cm, height=4.2cm,
    xlabel={$V$ (V)}, ylabel={$I$ (A)},
    xmin=0, xmax=5, ymin=0, ymax=5,
    grid=major, grid style={gray!15},
    xtick={0,1,2,3,4,5}, ytick={0,1,2,3,4,5},
    title={Filament Bulb (non-ohmic)}
  ]
    \addplot[codexred, very thick, smooth] coordinates {
      (0,0) (0.25,0.72) (0.5,1.2) (1,1.88) (1.5,2.38)
      (2,2.78) (2.5,3.1) (3,3.38) (4,3.86) (5,4.3)
    };
  \end{axis}

  % Diode
  \begin{axis}[
    name=p3,
    at={(p2.right of south east)},
    anchor=left of south west,
    xshift=0.5cm,
    width=4.0cm, height=4.2cm,
    xlabel={$V$ (V)}, ylabel={$I$ (A)},
    xmin=-2, xmax=3, ymin=-0.3, ymax=5,
    grid=major, grid style={gray!15},
    xtick={-2,-1,0,1,2,3},
    xticklabel style={font=\scriptsize},
    ytick={0,1,2,3,4,5},
    title={Diode (non-ohmic)}
  ]
    % Schematic forward-bias rise; the plotted range ends inside the axes.
    \addplot[codexgreen, very thick, domain=0:1.5,
      samples=60]{0.5*(exp(1.5*x)-1)};
    % Reverse leakage is negligible on this scale; breakdown is omitted.
    \addplot[codexgreen, very thick, domain=-2:0]{-0.00001};
    \node[codexgreen, font=\tiny, align=center,
      text width=0.9cm, inner sep=0.5pt, anchor=east]
      at (axis cs:2.8,4.45) {Forward\\bias};
  \end{axis}
  \node[codexblue, font=\footnotesize, align=center,
    text width=2.8cm, inner sep=1pt, anchor=north east, yshift=-1.1cm]
    at (p1.south east) {Straight line;\\$R$ constant};
  \node[codexred, font=\footnotesize, align=center,
    text width=2.8cm, inner sep=1pt, anchor=north east, yshift=-1.1cm]
    at (p2.south east) {$R$ rises as bulb\\warms};
  \node[codexred, font=\footnotesize, align=center,
    text width=3.6cm, inner sep=1pt, anchor=north east, yshift=-1.1cm]
    at (p3.south east) {Reverse leakage\\negligible on this\\scale; breakdown\\omitted};
\end{tikzpicture}
\end{center}
\end{diagbox}

% ============================================================
\section{Resistivity}
% ============================================================

\begin{definitionbox}[Resistivity]
The \textbf{resistivity} $\rho$ of a material is
its inherent resistance per unit length per unit
cross-sectional area:
\[
  R = \frac{\rho L}{A}
\]
where $L$ is the length of the conductor and $A$
is its cross-sectional area.

SI unit: $\Omega$\,m (ohm-metre).

Resistivity is a \textbf{material property};
resistance is a \textbf{property of a specific
component}. Two wires of different lengths made
from the same material have different resistances
but the same resistivity.
\end{definitionbox}

\begin{table}[H]
\centering
\caption{Resistivities of Common Materials at 20°C}
\label{tab:resistivity}
\begin{tabular}{lll}
\toprule
\textbf{Material} & \textbf{Resistivity ($\Omega$\,m)} &
\textbf{Type}\\
\midrule
Silver & $1.59\times10^{-8}$ & Conductor\\
Copper & $1.72\times10^{-8}$ & Conductor\\
Gold & $2.44\times10^{-8}$ & Conductor\\
Aluminium & $2.82\times10^{-8}$ & Conductor\\
Iron & $1.0\times10^{-7}$ & Conductor\\
Nichrome & $1.1\times10^{-6}$ & Alloy (heating elements)\\
Carbon (graphite) & $3.5\times10^{-5}$ & Semi-conductor\\
Silicon & $6.4\times10^{2}$ & Semiconductor\\
Glass & $10^{10}$--$10^{14}$ & Insulator\\
Rubber & $10^{13}$--$10^{15}$ & Insulator\\
\bottomrule
\end{tabular}
\end{table}

\begin{examplebox}[23.2]
\stepcounter{workedexample}
A nichrome heating element is $0.80\ \text{m}$ long
and has a diameter of $0.50\ \text{mm}$.
($\rho_{nichrome} = 1.1\times10^{-6}\ \Omega\text{m}$)\\
(a) Calculate the resistance of the element.\\
(b) It is connected to a $240\ \text{V}$ supply.
Find the current and power dissipated.\\
(c) A second identical element is connected in
parallel. Find the total resistance and power.
\end{examplebox}

\begin{solutionbox}
\textbf{(a)} $A = \pi r^2 = \pi(0.25\times10^{-3})^2
= 1.963\times10^{-7}\ \text{m}^2$:
\[
  R = \frac{\rho L}{A}
  = \frac{1.1\times10^{-6} \times 0.80}
    {1.963\times10^{-7}}
  = \frac{8.8\times10^{-7}}{1.963\times10^{-7}}
  = 4.48\ \Omega \approx 4.5\ \Omega
\]

\textbf{(b)} Current: $I = V/R = 240/4.48
= 53.6\ \text{A}$\\
Power: $P = VI = 240 \times 53.6 = 12\,864\ \text{W}$
$\approx 12.9\ \text{kW}$

\textbf{(c)} Two identical resistors in parallel:
$R_{total} = R/2 = 2.24\ \Omega$\\
Power: $P = V^2/R_{total} = 240^2/2.24
= 25\,714\ \text{W} = 25.7\ \text{kW}$

Parallel connection doubles the current and
doubles the power --- two independent paths for
current flow.
\end{solutionbox}

% ============================================================
\section{Resistance and Temperature}
% ============================================================

\begin{defbox}[Variation of Resistance with Temperature]
\textbf{Metals (conductors):}
Resistance \textit{increases} with temperature.
Higher temperature causes more vigorous lattice
vibrations, more frequent collisions between
electrons and ions, lower drift velocity.
\[
  R = R_0(1 + \alpha\Delta T)
\]
where $\alpha$ is the temperature coefficient of
resistance (K$^{-1}$).

\textbf{Semiconductors and thermistors (NTC type):}
Resistance \textit{decreases} with temperature.
Higher temperature frees more electrons as charge
carriers, greatly increasing conductivity.
(NTC = Negative Temperature Coefficient)

\textbf{LDR (Light Dependent Resistor):}
Resistance decreases as light intensity increases.
Light photons free more charge carriers.
\end{defbox}

\begin{diagbox}[Qualitative Normalized Resistance Trends]
\begin{center}
\begin{tikzpicture}
\begin{axis}[
  width=10cm, height=5cm,
  xlabel={Temperature $T$ (°C)},
  ylabel={Normalized resistance $R(T)/R(20^\circ\mathrm{C})$},
  xmin=0, xmax=120,
  ymin=0, ymax=2.5,
  grid=major, grid style={gray!15},
  ytick={0,0.5,1,1.5,2,2.5},
  xtick={0,20,40,60,80,100,120},
]
  % Each material is normalized to its own resistance at 20 degrees C.
  \addplot[codexblue, very thick, domain=0:120]
    {1.0 + 0.005*(x-20)};
  \addplot[codexred, very thick, domain=0:120,
    samples=80]
    {exp(-0.04*(x-20))};
\end{axis}
\end{tikzpicture}
\end{center}
{\footnotesize\centering
\textit{Each curve is normalized to its own resistance at $20^\circ\mathrm{C}$;
absolute resistance values are not comparable.}\par}
\smallskip
\noindent
\begin{minipage}[t]{0.48\linewidth}
  \footnotesize\raggedright
  \textcolor{codexblue}{\textbf{Metal:}} More lattice vibrations increase
  electron scattering, so resistance rises with temperature.
\end{minipage}\hfill
\begin{minipage}[t]{0.48\linewidth}
  \footnotesize\raggedright
  \textcolor{codexred}{\textbf{NTC thermistor:}} Heating frees more charge
  carriers, so resistance falls.
\end{minipage}
\end{diagbox}

% ============================================================
\section{Resistors in Series and Parallel}
% ============================================================

\begin{lawbox}[Resistors in Series]
Resistors in series carry the \textbf{same current}.
The total voltage is the sum of individual voltages.
\[
  V = V_1 + V_2 + V_3 + \ldots
\]
\[
  IR_{total} = IR_1 + IR_2 + IR_3 + \ldots
\]
\[
  \boxed{R_{series} = R_1 + R_2 + R_3 + \ldots}
\]
\end{lawbox}

\begin{lawbox}[Resistors in Parallel]
Resistors in parallel have the \textbf{same voltage}.
The total current is the sum of individual currents.
\[
  I = I_1 + I_2 + I_3 + \ldots
\]
\[
  \frac{V}{R_{total}} = \frac{V}{R_1}
  + \frac{V}{R_2} + \frac{V}{R_3} + \ldots
\]
\[
  \boxed{\frac{1}{R_{parallel}} = \frac{1}{R_1}
  + \frac{1}{R_2} + \frac{1}{R_3} + \ldots}
\]
For two resistors in parallel:
$R_p = \frac{R_1 R_2}{R_1 + R_2}$
\end{lawbox}

\begin{examplebox}[23.3]
\stepcounter{workedexample}
Find the total resistance and current from a
$24\ \text{V}$ supply for the network: $R_1 = 4\ \Omega$
in series with the parallel combination of
$R_2 = 6\ \Omega$ and $R_3 = 12\ \Omega$.
\end{examplebox}

\begin{solutionbox}
Parallel combination of $R_2$ and $R_3$:
\[
  R_{23} = \frac{R_2 R_3}{R_2 + R_3}
  = \frac{6 \times 12}{6 + 12}
  = \frac{72}{18} = 4\ \Omega
\]

Total resistance ($R_1$ in series with $R_{23}$):
\[
  R_{total} = R_1 + R_{23} = 4 + 4 = 8\ \Omega
\]

Total current:
\[
  I = \frac{V}{R_{total}} = \frac{24}{8}
  = 3.0\ \text{A}
\]

Voltage across parallel section:
$V_{23} = I \times R_{23} = 3.0 \times 4 = 12\ \text{V}$

Current through $R_2$: $I_2 = 12/6 = 2.0\ \text{A}$

Current through $R_3$: $I_3 = 12/12 = 1.0\ \text{A}$

Check: $I_2 + I_3 = 3.0\ \text{A} = I$ ✓
\end{solutionbox}

% ============================================================
\section{Kirchhoff's Laws}
% ============================================================

\begin{lawbox}[Kirchhoff's Current Law (KCL)]
\textbf{First Law (Junction Rule):}
The sum of currents entering any junction (node)
in a circuit equals the sum of currents leaving it.
\[
  \sum I_{in} = \sum I_{out}
\]
This is a consequence of conservation of charge
--- charge cannot accumulate at a junction.
\end{lawbox}

\begin{lawbox}[Kirchhoff's Voltage Law (KVL)]
\textbf{Second Law (Loop Rule):}
The algebraic sum of potential differences around
any closed loop in a circuit is zero.
\[
  \sum V = 0 \quad \text{(around any closed loop)}
\]
This is a consequence of conservation of energy
--- the work done going around a complete loop is
zero (you return to the same potential).
\end{lawbox}

\begin{diagbox}[Applying Kirchhoff's Laws]
\begin{center}
\begin{tikzpicture}[scale=0.95, font=\small]

  % Circuit nodes
  \coordinate (A) at (0,4);
  \coordinate (B) at (5,4);
  \coordinate (C) at (5,0);
  \coordinate (D) at (0,0);
  \coordinate (J) at (2.5,4);

  % Wires
  \draw[very thick] (A) -- (B) -- (C) -- (D) -- (A);
  \draw[very thick] (A) -- ++(2.5,0) -- ++(0,-4);

  % Batteries
  % E1 (left branch, top to bottom)
  \draw[very thick] (0,2.5) -- (0,2.8);
  % NOTE: use ASCII '-' for the coordinate; Unicode minus (−) breaks pgf math under pdfLaTeX
  \draw[very thick, codexred] (-0.3,2.8) -- (0.3,2.8);
  \draw[very thick] (0,2.8) -- (0,3.1);
  \draw[very thick, codexblue] (-0.5,3.1) -- (0.5,3.1);
  \draw[very thick] (0,3.1) -- (0,3.5);
  \node[left] at (-0.6,3.0) {$\mathcal{E}_1$};
  \node[right, font=\tiny] at (0.6,3.0) {EMF};

  % R1 on the outer branch before the central junction
  \fill[gray!30] (0.8,3.7) rectangle (1.8,4.3);
  \draw[thick] (0.8,3.7) rectangle (1.8,4.3);
  \node at (1.3,4.0) {$R_1$};

  % R2 (right branch)
  \fill[gray!30] (4.7,1.5) rectangle (5.3,3.0);
  \draw[thick] (4.7,1.5) rectangle (5.3,3.0);
  \node at (5.0,2.25) {$R_2$};

  % R3 (middle branch)
  \fill[gray!30] (2.2,1.5) rectangle (2.8,3.0);
  \draw[thick] (2.2,1.5) rectangle (2.8,3.0);
  \node at (2.5,2.25) {$R_3$};

  % E2 (middle branch; positive terminal at the bottom)
  \draw[very thick] (2.5,0.5) -- (2.5,0.8);
  \draw[very thick, codexblue] (2.0,0.8) -- (3.0,0.8);
  \draw[very thick] (2.5,0.8) -- (2.5,1.1);
  \draw[very thick, codexred] (2.2,1.1) -- (2.8,1.1);
  \node[right] at (3.15,0.95) {$\mathcal{E}_2$};

  % Current labels
  \draw[->, very thick, codexblue]
    (1.9,4.15) -- (2.4,4.15)
    node[above]{$I_1$};
  \draw[->, very thick, codexred]
    (B) ++ (0,-0.5) -- ++(0,-0.8)
    node[right]{$I_2$};
  \draw[->, very thick, codexgreen]
    (2.5,3.5) -- (2.5,3.0)
    node[right]{$I_3$};

  % KCL annotation
  \fill[codexgold] (J) circle (5pt);
  \node[anchor=west, codexgold, font=\tiny, align=left] at (2.8,4.55)
    {Junction A:\\$I_1 = I_2 + I_3$};

  % KVL equation callout; the main circuit shows the actual outer-loop path.
  \node[codexgold, font=\tiny, fill=white, inner sep=1pt,
        align=center, text width=3.2cm] at (2.5,-0.45)
    {Outer loop:\\$\mathcal{E}_1 - I_1R_1 - I_2R_2 = 0$};

\end{tikzpicture}
\end{center}
\end{diagbox}

\begin{examplebox}[23.4]
\stepcounter{workedexample}
In the circuit, $\mathcal{E}_1 = 12\ \text{V}$,
$\mathcal{E}_2 = 6\ \text{V}$,
$R_1 = 2\ \Omega$, $R_2 = 4\ \Omega$,
$R_3 = 6\ \Omega$. Find the three currents using
Kirchhoff's Laws.
\end{examplebox}

\begin{solutionbox}
Let currents $I_1$, $I_2$, $I_3$ flow as shown.

\textbf{KCL at junction A:}
\[
  I_1 = I_2 + I_3 \quad \cdots(1)
\]

\textbf{KVL, outer loop ($\mathcal{E}_1$, $R_1$, $R_2$):}
\[
  12 - 2I_1 - 4I_2 = 0
  \implies 2I_1 + 4I_2 = 12
  \implies I_1 + 2I_2 = 6 \quad\cdots(2)
\]

\textbf{KVL, inner loop ($R_3$, $R_2$):}
\[
  6I_3 - 4I_2 - 6 = 0
  \implies 6I_3 - 4I_2 = 6
  \implies 3I_3 - 2I_2 = 3 \quad\cdots(3)
\]

From (1): $I_1 = I_2 + I_3$. Substituting into (2):
\[
  (I_2 + I_3) + 2I_2 = 6
  \implies 3I_2 + I_3 = 6 \quad\cdots(4)
\]

From (3): $I_3 = (3 + 2I_2)/3$. Substituting into (4):
\[
  3I_2 + \frac{3 + 2I_2}{3} = 6
  \implies 9I_2 + 3 + 2I_2 = 18
  \implies 11I_2 = 15
  \implies I_2 = \frac{15}{11} \approx 1.36\ \text{A}
\]
\[
  I_3 = \frac{3 + 2(1.36)}{3}
  \approx 1.91\ \text{A}
\]
\[
  I_1 = I_2 + I_3 = 1.36 + 1.91 = 3.27\ \text{A}
\]
\end{solutionbox}

% ============================================================
\section{EMF and Internal Resistance}
% ============================================================

\begin{definitionbox}[EMF and Internal Resistance]
The \textbf{electromotive force} (EMF) $\mathcal{E}$
of a source is the energy supplied per unit charge
by the source (battery, generator, solar cell):
\[
  \mathcal{E} = \frac{W}{Q}
\]
SI unit: volt (V). EMF is \textbf{not} a force ---
it is a potential difference, but one that drives
the circuit.

Every real source has an \textbf{internal resistance}
$r$ (due to the resistance of its own materials).
For a source of EMF $\mathcal{E}$ and internal
resistance $r$ driving current $I$ through external
resistance $R$:
\[
  \mathcal{E} = IR + Ir = I(R + r)
\]
\[
  V_{terminal} = \mathcal{E} - Ir
\]
where $Ir$ is the \textbf{lost volts} --- potential
difference ``wasted'' inside the source.
\end{definitionbox}

\begin{diagbox}[Battery with Internal Resistance]
\begin{center}
\begin{tikzpicture}[scale=1.0, font=\small]

  % The source body contains the ideal EMF and its series internal resistance.
  \draw[thick, dashed, gray]
    (0,-0.5) rectangle (3.5,1.5);
  \node[gray, font=\scriptsize, fill=white, inner sep=1pt]
    at (0.9,1.2) {Real source};

  % External circuit: load R is in series with the real source.
  \draw[very thick] (0,0.5) -- (-1.0,0.5)
    -- (-1.0,2.5);
  \draw[very thick] (3.5,0.5) -- (4.5,0.5)
    -- (4.5,2.5);
  \draw[very thick] (-1.0,2.5) -- (0.5,2.5);
  \draw[very thick] (3.0,2.5) -- (4.5,2.5);

  % External resistor
  \fill[codexblue!20] (0.5,2.2) rectangle (3.0,2.8);
  \draw[thick] (0.5,2.2) rectangle (3.0,2.8);
  \node[above, font=\small] at (1.75,2.8) {External load};
  \node at (1.75,2.5) {$R$};

  % Internal branch: terminal, r, ideal EMF, return terminal.
  \draw[very thick] (0,0.5) -- (0.3,0.5);
  \fill[gray!30] (0.3,0.3) rectangle (1.6,0.7);
  \draw[thick] (0.3,0.3) rectangle (1.6,0.7);
  \node[font=\tiny] at (0.95,0.5) {$r$};
  \draw[very thick] (1.6,0.5) -- (2.0,0.5);
  \draw[very thick, codexblue] (2.0,0.15) -- (2.0,0.85);
  \draw[very thick, codexred] (2.3,0.3) -- (2.3,0.7);
  \draw[very thick] (2.3,0.5) -- (3.5,0.5);
  \node[above, font=\tiny] at (2.15,0.9) {$\mathcal{E}$};
  \node[codexblue, font=\tiny] at (1.9,0.9) {$+$};

  % Current arrow
  \draw[->, very thick, codexgold]
    (3.7,2.5) -- (4.35,2.5)
    node[above, font=\small]{$I$};

  % Voltage labels
  \draw[<->, codexblue, thick]
    (-1.5,0.5) -- (-1.5,2.5);
  \node[left, codexblue] at (-1.6,1.5)
    {$V_{terminal}$\\$= \mathcal{E} - Ir$};

  % Lost-voltage drop across the internal resistor.
  \draw[<->, codexred, thick] (0.3,0.12) -- (1.6,0.12);
  \node[below, codexred, font=\tiny] at (0.95,0.12)
    {$Ir$ (lost volts)};

  % Equation
  \node[align=center, font=\small]
    at (1.75,-1.15)
    {$\mathcal{E} = V_{terminal} + Ir
    = I(R+r)$};

\end{tikzpicture}
\end{center}
\end{diagbox}

\begin{examplebox}[23.5]
\stepcounter{workedexample}
A battery of EMF $9.0\ \text{V}$ and internal
resistance $0.5\ \Omega$ is connected to an external
resistor of $8.5\ \Omega$.\\
(a) Find the current in the circuit.\\
(b) Find the terminal voltage of the battery.\\
(c) Find the power delivered to the external resistor.\\
(d) What is the maximum power the battery could
deliver (short circuit current), and why is this
dangerous?
\end{examplebox}

\begin{solutionbox}
\textbf{(a)} Current:
\[
  I = \frac{\mathcal{E}}{R+r}
  = \frac{9.0}{8.5+0.5} = \frac{9.0}{9.0}
  = 1.0\ \text{A}
\]

\textbf{(b)} Terminal voltage:
\[
  V_T = \mathcal{E} - Ir = 9.0 - 1.0(0.5)
  = 8.5\ \text{V}
\]
or: $V_T = IR = 1.0 \times 8.5 = 8.5\ \text{V}$ ✓

\textbf{(c)} Power to external resistor:
\[
  P = I^2R = (1.0)^2 \times 8.5 = 8.5\ \text{W}
\]

\textbf{(d)} Short circuit ($R = 0$):
\[
  I_{sc} = \frac{\mathcal{E}}{r}
  = \frac{9.0}{0.5} = 18\ \text{A}
\]
Power dissipated inside the battery:
$P_{internal} = I^2 r = (18)^2 \times 0.5 = 162\ \text{W}$

This enormous power dissipation inside the battery
causes rapid heating, potentially causing the
battery to explode or catch fire. Short circuits
are genuinely dangerous.
\end{solutionbox}

% ============================================================
\section{The Potential Divider}
% ============================================================

\begin{definitionbox}[Potential Divider]
A \textbf{potential divider} (voltage divider) is
two (or more) resistors connected in series across
a supply, used to produce an output voltage that
is a fraction of the supply voltage.

For two resistors $R_1$ and $R_2$ in series across
supply $V_s$, the output voltage across $R_2$ is:
\[
  V_{out} = V_s \times \frac{R_2}{R_1 + R_2}
\]
\end{definitionbox}

\begin{diagbox}[Potential Divider Circuit]
\begin{center}
\begin{tikzpicture}[scale=1.0, font=\small]

  % Series divider string and supply rails.
  \draw[very thick] (0,5.0) -- (4.0,5.0);
  \draw[very thick] (0,0.0) -- (4.0,0.0);
  \draw[very thick] (0,5.0) -- (0,4.5);
  \draw[very thick] (0,3.5) -- (0,3.0);
  \draw[very thick] (0,3.0) -- (0,2.0);
  \draw[very thick] (0,1.0) -- (0,0.0);

  % R1 and R2 are actually in series between the supply rails.
  \fill[gray!30] (-0.3,3.5) rectangle (0.3,4.5);
  \draw[thick] (-0.3,3.5) rectangle (0.3,4.5);
  \node[left] at (-0.5,4.0) {$R_1$};
  \fill[gray!30] (-0.3,1.0) rectangle (0.3,2.0);
  \draw[thick] (-0.3,1.0) rectangle (0.3,2.0);
  \node[left] at (-0.5,1.5) {$R_2$};
  \fill[black] (0,3.0) circle (4pt);
  \fill[black] (0,0.0) circle (4pt);

  % Connected source in the return rail (long plate is positive).
  \draw[very thick] (4.0,0.0) -- (4.0,2.7);
  \draw[very thick] (4.0,3.0) -- (4.0,5.0);
  \draw[very thick, codexblue] (3.55,3.0) -- (4.45,3.0);
  \draw[very thick, codexred] (3.75,2.7) -- (4.25,2.7);
  \node[left] at (3.45,2.85) {$V_s$};
  \node[right, codexblue, font=\tiny] at (4.45,3.0) {$+$};
  \node[right, codexred, font=\tiny] at (4.25,2.7) {$-$};

  % Open voltage-measurement taps across R2: no wire shorts the output.
  \draw[thick] (0,3.0) -- (1.4,3.0);
  \draw[thick] (0,0.0) -- (1.4,0.0);
  \draw[<->, codexblue, thick] (1.8,0.0) -- (1.8,3.0);
  \node[right, codexblue] at (1.9,1.5) {$V_{out}$};

  % Thermistor variant annotation
  \node[align=center, font=\footnotesize, text width=2.8cm]
    at (6.1,2.8)
    {If $R_2$ is a thermistor:\
    $V_{out}$ varies with temperature\\[0.3cm]
    If $R_2$ is an LDR:\
    $V_{out}$ varies with light intensity};

  % Formula
  \node[codexblue, font=\normalsize] at (6.1,0.55)
    {$V_{out} = V_s\dfrac{R_2}{R_1+R_2}$};

\end{tikzpicture}
\end{center}
\end{diagbox}

\begin{examplebox}[23.6]
\stepcounter{workedexample}
A potential divider consists of a $10\ \text{k}\Omega$
resistor in series with an NTC thermistor. The supply
is $5.0\ \text{V}$. The thermistor has resistance
$20\ \text{k}\Omega$ at $20°\text{C}$ and
$2\ \text{k}\Omega$ at $80°\text{C}$.
The output is taken across the thermistor.\\
(a) Find $V_{out}$ at $20°\text{C}$.\\
(b) Find $V_{out}$ at $80°\text{C}$.\\
(c) Explain how this circuit could be used as a
temperature-sensing switch.
\end{examplebox}

\begin{solutionbox}
\textbf{(a)} At $20°\text{C}$ ($R_T = 20\ \text{k}\Omega$):
\[
  V_{out} = 5.0 \times \frac{20}{10+20}
  = 5.0 \times \frac{20}{30} = 3.33\ \text{V}
\]

\textbf{(b)} At $80°\text{C}$ ($R_T = 2\ \text{k}\Omega$):
\[
  V_{out} = 5.0 \times \frac{2}{10+2}
  = 5.0 \times \frac{2}{12} = 0.833\ \text{V}
\]

\textbf{(c)} $V_{out}$ decreases as temperature
rises (because the thermistor resistance falls,
taking less of the supply voltage). This output
can be connected to a comparator or transistor
circuit that switches a relay at a set voltage
threshold. For example: if $V_{out}$ falls below
$1.0\ \text{V}$ (corresponding to a temperature
above about $75°\text{C}$), a fan or cooling
system switches on. This is exactly how temperature
control circuits work in computers, refrigerators,
and air conditioning units.
\end{solutionbox}

% ============================================================
\section{The Wheatstone Bridge}
% ============================================================

\begin{definitionbox}[Wheatstone Bridge]
A \textbf{Wheatstone bridge} is a circuit for
measuring an unknown resistance by comparison with
known resistances. It consists of four resistors
$P$, $Q$, $R$, $S$ arranged in a diamond with a
galvanometer $G$ across the middle.

The bridge is \textbf{balanced} (zero galvanometer
current) when:
\[
  \frac{P}{Q} = \frac{R}{S}
  \implies PS = QR
\]
At balance, the unknown resistance $S$ can be found:
\[
  S = \frac{QR}{P}
\]
\end{definitionbox}

\begin{diagbox}[Wheatstone Bridge Circuit]
\begin{center}
\begin{tikzpicture}[scale=0.9, font=\small]

  % Four bridge junctions.
  \coordinate (A) at (0,3);
  \coordinate (B) at (3,6);
  \coordinate (C) at (6,3);
  \coordinate (D) at (3,0);

  % Each resistor interrupts its own diamond arm.
  \draw[thick] (A) -- (1.2,4.2);
  \node[draw, fill=gray!25, minimum width=0.8cm,
        minimum height=0.35cm, rotate=45] at (1.5,4.5) {$P$};
  \draw[thick] (1.8,4.8) -- (B);

  \draw[thick] (B) -- (4.2,4.8);
  \node[draw, fill=gray!25, minimum width=0.8cm,
        minimum height=0.35cm, rotate=-45] at (4.5,4.5) {$Q$};
  \draw[thick] (4.8,4.2) -- (C);

  \draw[thick] (A) -- (1.2,1.8);
  \node[draw, fill=gray!25, minimum width=0.8cm,
        minimum height=0.35cm, rotate=-45] at (1.5,1.5) {$R$};
  \draw[thick] (1.8,1.2) -- (D);

  \draw[thick] (D) -- (4.2,1.2);
  \node[draw, fill=gray!25, minimum width=0.8cm,
        minimum height=0.35cm, rotate=45] at (4.5,1.5) {$S$};
  \draw[thick] (4.8,1.8) -- (C);

  % Galvanometer connects the two middle junctions B and D.
  \draw[thick] (B) -- (3,3.4);
  \fill[white] (2.6,2.6) rectangle (3.4,3.4);
  \draw[thick] (2.6,2.6) rectangle (3.4,3.4);
  \node at (3,3) {$G$};
  \draw[thick] (3,2.6) -- (D);
  \node[right, font=\tiny] at (3.5,3) {Galvanometer};

  % Supply is connected across A and C via an external branch above B.
  \draw[thick] (A) -- (-1,3) -- (-1,7.3) -- (2.7,7.3);
  \draw[very thick, codexblue] (2.7,6.95) -- (2.7,7.65);
  \draw[very thick, codexred] (3.05,7.1) -- (3.05,7.5);
  \draw[thick] (3.05,7.3) -- (7,7.3) -- (7,3) -- (C);
  \node[above, font=\small] at (2.9,7.7) {$V_s$};
  \node[codexblue, font=\tiny] at (2.55,7.7) {$+$};

  % Junction markers and labels.
  \foreach \coord in {(A),(B),(C),(D)} {
    \fill[black] \coord circle (4pt);
  }
  \node[below left] at (A) {A};
  \node[above] at (B) {B};
  \node[below right] at (C) {C};
  \node[right] at (3.2,0) {D};

  % Balance condition.
  \node[align=center, font=\small] at (3,-1.0)
    {At balance ($I_G = 0$):\\
    $\dfrac{P}{Q} = \dfrac{R}{S}$};

\end{tikzpicture}
\end{center}
\end{diagbox}

% ============================================================
\section{The Potentiometer}
% ============================================================

\begin{definitionbox}[Potentiometer]
A \textbf{potentiometer} is a device for measuring
EMF or potential difference by comparison, using a
null method (no current drawn from the source being
measured).

A long resistance wire of length $L$ is connected
to a driver cell. The potential per unit length along
the wire is constant. A sliding contact (jockey) is
moved along the wire until the galvanometer reads
zero (null point). At the null point:
\[
  \frac{\mathcal{E}_1}{\mathcal{E}_2}
  = \frac{\ell_1}{\ell_2}
\]
where $\ell_1$ and $\ell_2$ are the balancing lengths
for EMFs $\mathcal{E}_1$ and $\mathcal{E}_2$.

\textbf{Advantage over voltmeter:} At the null point,
no current is drawn from the cell being measured, so
its terminal voltage equals its EMF.
\end{definitionbox}

% ============================================================
\section{Chapter Summary}
% ============================================================

\begin{summarybox}
\textbf{Current:} $I = \Delta Q/\Delta t$;
$I = nqAv_d$ (drift velocity)

\textbf{Ohm's Law:} $V = IR$ (ohmic, constant $T$)

\textbf{Resistivity:} $R = \rho L/A$;
metals: $\rho$ increases with $T$;
NTC thermistors: $\rho$ decreases with $T$.

\textbf{Series:} $R_s = R_1+R_2+\ldots$
(same current; voltages add)

\textbf{Parallel:}
$\frac{1}{R_p} = \frac{1}{R_1}+\frac{1}{R_2}+\ldots$
(same voltage; currents add)

\textbf{Kirchhoff's Laws:}
KCL: $\sum I_{in} = \sum I_{out}$ (charge conservation)
KVL: $\sum V = 0$ (energy conservation)

\textbf{EMF and internal resistance:}
$\mathcal{E} = I(R+r)$;
$V_{terminal} = \mathcal{E} - Ir$

\textbf{Potential divider:}
$V_{out} = V_s R_2/(R_1+R_2)$

\textbf{Wheatstone bridge balance:}
$P/Q = R/S$

\textbf{Potentiometer:} null method;
$\mathcal{E}_1/\mathcal{E}_2 = \ell_1/\ell_2$
\end{summarybox}

% ============================================================
\newpage
\begin{exercisebox}[23]

\subsection*{Section A: Multiple Choice}

\textbf{A1.} A charge of $120\ \text{C}$ flows
through a resistor in $4.0\ \text{min}$. The
current is:

\mcoption{A}{$0.5\ \text{A}$}
\mcoption{B}{$2.0\ \text{A}$}
\mcoption{C}{$30\ \text{A}$}
\mcoption{D}{$480\ \text{A}$}
\ansref{Ch23-A1}

\vspace{0.4cm}

\textbf{A2.} Three resistors of $6\ \Omega$ each
are connected: one in series with the parallel
combination of the other two. The total resistance is:

\mcoption{A}{$18\ \Omega$}
\mcoption{B}{$9\ \Omega$}
\mcoption{C}{$3\ \Omega$}
\mcoption{D}{$2\ \Omega$}
\ansref{Ch23-A2}

\vspace{0.4cm}

\textbf{A3.} A cell of EMF $6.0\ \text{V}$ and
internal resistance $1.0\ \Omega$ drives a current
of $2.0\ \text{A}$. The terminal voltage is:

\mcoption{A}{$4.0\ \text{V}$}
\mcoption{B}{$5.0\ \text{V}$}
\mcoption{C}{$6.0\ \text{V}$}
\mcoption{D}{$8.0\ \text{V}$}
\ansref{Ch23-A3}

\vspace{0.4cm}

\textbf{A4.} Kirchhoff's Voltage Law is a
consequence of:

\mcoption{A}{Conservation of charge}
\mcoption{B}{Conservation of energy}
\mcoption{C}{Newton's Third Law}
\mcoption{D}{Faraday's Law}
\ansref{Ch23-A4}

\vspace{0.4cm}

\textbf{A5.} In a potential divider, $R_1 = 3\ \Omega$
and $R_2 = 9\ \Omega$ across a $12\ \text{V}$ supply.
The output voltage across $R_2$ is:

\mcoption{A}{$3\ \text{V}$}
\mcoption{B}{$6\ \text{V}$}
\mcoption{C}{$9\ \text{V}$}
\mcoption{D}{$12\ \text{V}$}
\ansref{Ch23-A5}

\vspace{0.4cm}

\textbf{A6.} The resistance of a metal wire
increases when:

\mcoption{A}{Its length decreases}
\mcoption{B}{Its cross-sectional area increases}
\mcoption{C}{Its temperature increases}
\mcoption{D}{Its temperature decreases}
\ansref{Ch23-A6}

\vspace{0.4cm}

\textbf{A7.} A Wheatstone bridge is balanced when
$P = 10\ \Omega$, $Q = 5\ \Omega$, $R = 8\ \Omega$.
The value of the unknown resistance $S$ is:

\mcoption{A}{$4\ \Omega$}
\mcoption{B}{$8\ \Omega$}
\mcoption{C}{$16\ \Omega$}
\mcoption{D}{$40\ \Omega$}
\ansref{Ch23-A7}

\vspace{0.4cm}

\textbf{A8.} An NTC thermistor in a circuit:

\mcoption{A}{Has resistance that increases with
light intensity}
\mcoption{B}{Has resistance that increases with
temperature}
\mcoption{C}{Has resistance that decreases with
temperature}
\mcoption{D}{Obeys Ohm's Law at all temperatures}
\ansref{Ch23-A8}

\vspace{0.4cm}

\textbf{A9.} The drift velocity of electrons
carrying $3\ \text{A}$ in a $2\ \text{mm}^2$ copper
wire ($n = 8.5\times10^{28}\ \text{m}^{-3}$,
$e = 1.6\times10^{-19}\ \text{C}$) is approximately:

\mcoption{A}{$1.1\times10^{-4}\ \text{m\,s}^{-1}$}
\mcoption{B}{$1.1\times10^{-3}\ \text{m\,s}^{-1}$}
\mcoption{C}{$3.0\times10^{-4}\ \text{m\,s}^{-1}$}
\mcoption{D}{$3.0\ \text{m\,s}^{-1}$}
\ansref{Ch23-A9}

\vspace{0.4cm}

\textbf{A10.} The potentiometer is preferred over
a voltmeter for measuring EMF because:

\mcoption{A}{The potentiometer is cheaper}
\mcoption{B}{At balance, no current is drawn from
the cell, so the terminal voltage equals the EMF}
\mcoption{C}{The potentiometer is more portable}
\mcoption{D}{Voltmeters cannot measure EMF at all}
\ansref{Ch23-A10}

\vspace{0.6cm}

\subsection*{Section B: Theory and Calculations}

\textbf{B1.} A copper wire of diameter $1.0\ \text{mm}$
and length $2.0\ \text{m}$ is connected to a
$3.0\ \text{V}$ battery.
($\rho_{Cu} = 1.72\times10^{-8}\ \Omega\text{m}$,
$n = 8.5\times10^{28}\ \text{m}^{-3}$,
$e = 1.6\times10^{-19}\ \text{C}$)\\[4pt]
(a) Calculate the resistance of the wire.\\
(b) Calculate the current flowing.\\
(c) Calculate the drift velocity of electrons.\\
(d) Comment on the comparison between drift velocity
and the speed at which the circuit switches on.\\
(e) The wire is replaced with one of the same length
and material but half the diameter. What is the new
resistance and current?
\hfill\ansref{Ch23-B1}

\vspace{0.4cm}

\textbf{B2.} Apply Kirchhoff's Laws to the following
circuit: Battery $\mathcal{E}_1 = 10\ \text{V}$
(internal resistance $0.5\ \Omega$) is connected
in a loop with $R_1 = 3\ \Omega$. A second battery
$\mathcal{E}_2 = 4\ \text{V}$ (internal resistance
$1.0\ \Omega$) and $R_2 = 2\ \Omega$ form a branch
connected in parallel with $R_1$.
A third branch with $R_3 = 6\ \Omega$ is also
in parallel with $R_1$.\\[4pt]
(a) Label the three branch currents $I_1$, $I_2$,
$I_3$ and write KCL at the junction.\\
(b) Write two independent KVL equations.\\
(c) Solve for the three currents.\\
(d) Verify your answers using power balance:
total power from sources = total power in resistors.
\hfill\ansref{Ch23-B2}

\vspace{0.4cm}

\textbf{B3.} A battery of EMF $12\ \text{V}$ is
connected to the following network:
$R_1 = 2\ \Omega$ in series with the parallel
combination of $R_2 = 6\ \Omega$ and $R_3 = 3\ \Omega$,
and then in series with $R_4 = 1\ \Omega$.
The battery has internal resistance $r = 0.5\ \Omega$.\\[4pt]
(a) Find the total circuit resistance.\\
(b) Find the current from the battery.\\
(c) Find the terminal voltage.\\
(d) Find the voltage across $R_2$.\\
(e) Find the current through $R_3$.\\
(f) Find the power dissipated in $R_4$.
\hfill\ansref{Ch23-B3}

\begin{center}
\begin{tikzpicture}[scale=0.82,transform shape,font=\small,>=Stealth]
  % Source and return path
  \draw[thick] (0.65,0.4) -- (0.65,1.35);
  \draw[thick] (0.65,2.15) -- (0.65,3.1);
  \node[draw,circle,minimum size=8mm,inner sep=1pt] at (0.65,1.75) {$12\,\text{V}$};
  \draw[thick] (0.65,0.4) -- (7.35,0.4) -- (7.35,3.1);
  % Series internal resistance and R1
  \draw[thick] (0.65,3.1) -- (1.0,3.1);
  \draw[thick] (1.0,2.9) rectangle (1.8,3.3);
  \node at (1.4,3.1) {$r$};
  \node[above] at (1.4,3.3) {$0.5\,\Omega$};
  \draw[thick] (1.8,3.1) -- (2.05,3.1);
  \draw[thick] (2.05,2.9) rectangle (2.85,3.3);
  \node at (2.45,3.1) {$R_1$};
  \node[above] at (2.45,3.3) {$2\,\Omega$};
  \draw[thick] (2.85,3.1) -- (3.25,3.1);
  % Parallel R2 and R3 branches
  \draw[thick] (3.25,3.1) -- (3.55,3.1);
  \draw[thick] (3.55,2.9) rectangle (4.35,3.3);
  \node at (3.95,3.1) {$R_2$};
  \node[above] at (3.95,3.3) {$6\,\Omega$};
  \draw[thick] (4.35,3.1) -- (5.15,3.1);
  \draw[thick] (3.25,3.1) -- (3.25,1.35) -- (3.55,1.35);
  \draw[thick] (3.55,1.15) rectangle (4.35,1.55);
  \node at (3.95,1.35) {$R_3$};
  \node[below] at (3.95,1.15) {$3\,\Omega$};
  \draw[thick] (4.35,1.35) -- (5.15,1.35) -- (5.15,3.1);
  \draw[thick] (5.15,3.1) -- (5.45,3.1);
  % R4 remains in series after the branches recombine.
  \draw[thick] (5.45,2.9) rectangle (6.25,3.3);
  \node at (5.85,3.1) {$R_4$};
  \node[above] at (5.85,3.3) {$1\,\Omega$};
  \draw[thick] (6.25,3.1) -- (7.35,3.1);
  \node[left] at (0.65,3.1) {$+$};
\end{tikzpicture}
\end{center}

\vspace{0.4cm}

\textbf{B4.} In a Wheatstone bridge, $P = 100\ \Omega$,
$Q = 200\ \Omega$, and $R$ is a variable resistor.
$S$ is an unknown resistance (a strain gauge).
The driver cell is $6.0\ \text{V}$.\\[4pt]
(a) Find the value of $R$ needed to balance the
bridge when $S = 85\ \Omega$.\\
(b) When the strain gauge is stretched, its
resistance changes to $86.5\ \Omega$. The bridge
is no longer balanced. Calculate the new value of
$R$ needed to rebalance.\\
(c) The fractional change in resistance
$\Delta R/R = 0.0176$ corresponds to a strain of
$\varepsilon = 8.8\times10^{-5}$ (a tiny stretch).
Calculate the gauge factor $GF = (\Delta R/R)/\varepsilon$.\\
(d) Explain why a bridge circuit is better than a
simple voltage divider for detecting small resistance
changes.
\hfill\ansref{Ch23-B4}

\vspace{0.4cm}

\textbf{B5.} A student uses a potentiometer to
compare two cells. A driver cell of
$\mathcal{E}_d = 4.0\ \text{V}$ is connected across
a $100\ \text{cm}$ uniform resistance wire.
Cell 1 (standard cell, $\mathcal{E}_1 = 1.50\ \text{V}$)
balances at $\ell_1 = 37.5\ \text{cm}$.
Cell 2 (unknown EMF) balances at
$\ell_2 = 45.0\ \text{cm}$.\\[4pt]
(a) Calculate the potential gradient along the wire
(V per cm).\\
(b) Calculate the EMF of cell 2.\\
(c) Cell 2 is now connected with a resistance of
$5.0\ \Omega$ in series (a load), and the balancing
length changes to $35.0\ \text{cm}$.
Find the terminal voltage of cell 2 under this load.\\
(d) Calculate the internal resistance of cell 2.\\
(e) Explain why the potentiometer gives a more
accurate value of EMF than a voltmeter would.
\hfill\ansref{Ch23-B5}

\begin{center}
\begin{tikzpicture}[scale=0.9,transform shape,font=\small,>=Stealth]
  % The jockey is shown at 40% of AB as a representative movable null point.
  \node[font=\bfseries] at (2.4,3.78) {Potentiometer};
  \draw[thick] (0.3,2.65) -- (2.0,2.65) -- (4.55,2.65);
  \draw[thick] (0.3,2.65) -- (0.3,0.55) -- (1.9,0.55);
  \node[draw,circle,minimum size=8mm,inner sep=1pt] at (2.3,0.55) {$4.0\,\text{V}$};
  \draw[thick] (2.7,0.55) -- (4.55,0.55) -- (4.55,2.65);
  \node[fill=white,inner sep=0.6pt,font=\footnotesize] at (1.55,0.85) {$+$};
  \node[fill=white,inner sep=0.6pt,font=\footnotesize] at (3.05,0.85) {$-$};
  \node[below left] at (0.3,2.65) {$A$};
  \node[below right] at (4.55,2.65) {$B$};
  \node[font=\footnotesize] at (2.4,3.26)
    {uniform wire AB, $100\,\text{cm}$};
  \node[below] at (2.3,0.05) {driver cell};
  % J lies at (2.0-0.3)/(4.55-0.3)=40% of AB; the arrows show it can slide.
  \fill (2.0,2.65) circle (1.6pt);
  \node[below left,font=\scriptsize] at (2.0,2.65) {$J$};
  \draw[thin] (2.0,2.65) -- (2.0,2.95);
  \draw[<->,thin] (1.72,2.95) -- (2.28,2.95);
  \node[font=\tiny,align=center] at (3.8,2.0)
    {movable J /\\null point};
  % Test cell and galvanometer connect A to the jockey.
  \draw[thick] (0.3,2.65) -- (0.3,1.55) -- (0.55,1.55);
  \node[draw,minimum width=1.2cm,minimum height=0.5cm,inner sep=1pt] at (1.15,1.55) {$E_1/E_2$};
  \node[above] at (0.55,1.82) {$+$};
  \node[above] at (1.75,1.82) {$-$};
  \draw[thick] (1.75,1.55) -- (2.4,1.55);
  \node[draw,circle,minimum size=6mm,inner sep=1pt] at (2.7,1.55) {$G$};
  \draw[thick] (3.0,1.55) -- (2.0,2.65);
  % Loaded cell: internal resistance and external load in one closed loop.
  \begin{scope}[shift={(5.5,0.15)}]
    \node[font=\bfseries] at (1.7,3.65) {Cell under load};
    \draw[thick] (0.4,0.5) -- (0.4,1.05);
    \node[draw,circle,minimum size=8mm,inner sep=1pt] at (0.4,1.45) {$E_2$};
    \node[fill=white,inner sep=0.5pt,font=\footnotesize] at (0.88,1.83) {$+$};
    \node[fill=white,inner sep=0.5pt,font=\footnotesize] at (0.88,1.05) {$-$};
    \draw[thick] (0.4,1.85) -- (0.4,2.4) -- (1.0,2.4);
    \draw[thick] (1.0,2.2) rectangle (2.0,2.6);
    \node at (1.5,2.4) {$r_2$};
    \draw[thick] (2.0,2.4) -- (3.0,2.4) -- (3.0,1.8);
    \draw[thick] (2.7,1.8) rectangle (3.3,0.9);
    \node[right,font=\footnotesize] at (3.38,2.06) {$R_L=5\,\Omega$};
    \node[font=\footnotesize] at (3.4,1.8) {$+$};
    \node[font=\footnotesize] at (3.4,0.9) {$-$};
    \draw[<->,thin] (3.62,0.98) -- (3.62,1.72);
    \node[right,font=\footnotesize] at (3.72,1.35) {$V_{T2}$};
    \draw[thick] (3.0,0.9) -- (3.0,0.5) -- (0.4,0.5);
  \end{scope}
\end{tikzpicture}
\par\smallskip
\footnotesize Test cells are compared in turn; the $5\,\Omega$ load is only across cell 2 for the loaded measurement.
\end{center}

\end{exercisebox}
